% Page 008

\seite{8}

als die Quelle, aus der Sefferweich noch gegenwärtig seine\ob
Wasserleitung speist, in nächster Nähe ihren Ursprung hat.\ob
Ist der Sage der Alten hier Glauben beizumessen, so\ob
sollten sich in dieser Flur noch mehrere dieser Stein\ohb
mauern vorfinden, die der Bloßlegung noch harren.\ob
Ungefähr zehn Minuten von dieser Stelle entfernt\ob
wurde noch im Jahre 1886 am Abhange eines Berges\ob
(Örtel) nach der Abendseite ein römisches Grab geöffnet,\ob
schade daß es schon früher von Unberufenen dasselbe\ob
Schicksal erfahren; denn der Inhalt desselben, Asche\ob
und Thränenkrüglein lagen zerschlagen und zerstreut\ob
umher. Hat etwa eine römische Militärkolonie\ob
hier gelegen, welche zwischen den alten Trevieren\ob
und Belgern eingeschoben war! Ferner ist in der\ob
\margmark{|\\|\\|\\|\\|\\|}%
Nähe der Thüre der hiesigen Kapelle ein Stein mit\ob
lateinischer Schrift eingefügt, welche besagt, daß der\ohb
selbe früher einmal als Denkmal gedient, das\ob
ein Römer seiner verstorbenen Gattin in Sefferweich\ob
hatte errichten lassen. Ebenso wurde ein anderer\ob
Stein gefunden, welchen Sachkundige als Opfer\ohb
stein unserer alten heidnischen Vorfahren betrachten\ob
wissen wollen. Hier haben einmal Römer ge\ohb
herrscht.\ob
Einen ferneren Beweis für das hohe Alter Sefferweichs\ob
liegt in folgender Thatsache begründet: Im Pfarrarchiv\ob
zu Seffern befindet sich eine Pergamentrolle
\pend
